% !TeX root = ../main.tex
\section{Anomalies and Limitations}\label{sec:anomalies}
\rubric{Report quality}{10}{Clear plots, concise writing, honest anomalies and thoughtful limitations.}

\subsection{Unexpected results}
\begin{itemize}
  \item \textbf{Dot product never speeds up.} We expected a reduction to vectorize at $\approx8\times$
        like the other float32 kernels. Instead native dot stays at $\approx1\times$ across all sizes.
        Hypothesis: without \code{-ffast-math} the compiler may not reassociate floating-point
        addition, so it vectorizes the multiply but keeps a serial scalar sum. Tested: the
        disassembly shows \code{vmulps} on \code{ymm} followed by a chain of scalar \code{vaddss}
        (\cref{sec:vectorization}), confirming the hypothesis. We chose not to enable
        \code{-ffast-math} so that all variants remain bit-reproducible against the reference.
  \item \textbf{Native SAXPY speedup varies within L1.} It is $5.6\text{--}5.9\times$ at
        $N=\num{1024}$, \num{2048}, \num{3072} and \num{4096}, but $6.86\times$ at $N=\num{1536}$
        (locality sweep) and $6.32\times$ at $N=512$ (baseline). Scalar time is flat across these
        sizes and L1D misses stay below \num{0.0011} per element below $N=\num{4096}$; the peak is
        native needing \num{0.22} instead of \numrange{0.26}{0.28}~cycles per element. Per-call
        overhead cannot explain it, since it \emph{lowers} small-$N$ speedups
        (\cref{sec:alignment}). Untested hypothesis: the relative placement of \code{x} and
        \code{y} (4K aliasing, as for the stencil below), testable with the same distance sweep.
  \item \textbf{Stencil keeps a $2.6\text{--}2.8\times$ advantage in DRAM} while SAXPY and
        multiply fall to $1.88\text{--}2.02\times$ and $1.5\text{--}1.6\times$. Its native build is
        just as bandwidth-bound as theirs: it runs at \SI{101}{\percent} of the DRAM roof
        (\cref{sec:roofline}). Its scalar build, however, needs
        \num{5}~FLOP per element of serial \code{mulss}/\code{addss} work and reaches only
        \qtyrange{4.0}{5.0}{\giga\flop\per\second} at every size (\qty{4.0}{\giga\flop\per\second},
        \SI{37}{\percent} of the roof, in DRAM). The retained
        ratio is again scalar's instruction cost, as for SAXPY (\cref{sec:locality}), just larger.
  \item \textbf{Cache-resident stencil is slower than L3-resident stencil.} Native stencil delivers
        \qtyrange{16.6}{17.3}{\giga\flop\per\second} for $N=\num{1024}\ldots\num{16384}$ but
        \qtyrange{19.9}{20.7}{\giga\flop\per\second} for $N=\num{65536}\ldots\num{1048576}$, and
        \num{31.4} and \qty{27.4}{\giga\flop\per\second} at the isolated sizes $N=512$ and
        \num{32768}. The pattern reproduces exactly across separate runs; scalar stencil stays at \qtyrange{4.0}{5.0}{\giga\flop\per\second}. A cache
        effect would make the rate fall monotonically with $N$; this does not. Hypothesis (untested): 4K aliasing. The small \code{x} and \code{y}
        allocations sit at addresses whose difference makes a store to \code{y[i]} look like it
        conflicts with the load of \code{x[i+1]} in the low 12 address bits, which stalls the load;
        large allocations come from fresh pages with a different relative offset. It would be confirmed by a
        sweep of the \code{x}--\code{y} distance modulo \qty{4}{\kibi\byte} at fixed $N$, showing
        the slow rate only at small distances, together with the core's 4K-alias counter
        (\code{ld_blocks_partial.address_alias}). We ran neither.
  \item \textbf{Scalar and SIMD SAXPY do not converge in DRAM.} We expected one shared bandwidth
        plateau; native stays $1.88\text{--}2.02\times$ ahead because scalar SAXPY is still largely
        instruction-bound there (\num{2.34}~cycles per element at IPC \num{3.0}) while native sits on
        the roof (IPC \num{0.6}); the cycle ratio, $1.88$, is the measured speedup
        (\cref{sec:locality}).
  \item \textbf{Same code, different process, $N=\num{32768}$.} Scalar SAXPY is stable across
        processes elsewhere (\qtyrange{0.650}{0.653}{\nano\second} per element at all \num{9}
        locality sizes from $N=\num{65536}$ to \num{1048576}), but where the two arrays exactly fill
        the \qty{256}{\kibi\byte} L2 two measurements of the same machine code differ by
        \SI{1}{\percent} (scalar) and \SI{14}{\percent} (native); page placement is the untested
        hypothesis (\cref{sec:stride}).
  \item \textbf{Hardware gather loses to scalar.} The explicit \code{vgatherdps} kernel takes
        \qty{3.23}{\nano\second} per element at $N=\num{32768}$ versus \num{2.16} (native scalar
        loop) and \qty{2.02}{\nano\second} (scalar build), i.e.\ $1.5\text{--}1.6\times$
        \emph{slower}, and it is \SI{7}{\percent} slower in DRAM. Hypothesis: the Gather Data
        Sampling microcode mitigation, which this CPU runs (\cref{sec:stride}) and which is known
        to make \code{vgatherdps} much slower, plus the eight scalar
        stores that AVX2 cannot vectorize. Confirming it would need a boot without the mitigation,
        which requires root.
\end{itemize}

\subsection{Limitations}
\begin{itemize}
  \item \textbf{Core-PMU counters only.} The SIMD locality counters are core-PMU, user-space
        events of the benchmark thread (\cref{sec:ceilings-method}): cycles, instructions, generic LLC
        misses and L1D read misses. We did not count DRAM traffic for SIMD. The machine can do it:
        the companion Cache \& Memory report, run on the same machine the same day with
        \code{perf_event_paranoid}~$=0$, reads the memory-controller (\code{uncore_imc}) counters
        and finds IMC traffic matching its streaming estimate to within about \SI{1}{\percent}. TLB
        events were not recorded either. The generic LLC-miss event's treatment of hardware-prefetch
        requests is model-specific, so we use misses per element as a relative measure between
        builds rather than as exact DRAM traffic. Counting IMC bytes per element for scalar and
        native SAXPY would turn the roofline's traffic assumption into a measurement.
  \item \textbf{Fixed clock.} All reported experiments were run on 2026-10-21 with governor
        \sysGovernor{}, turbo \sysTurbo{} and SMT off; the cycle counter confirms
        \qtyrange{2.78}{2.79}{\giga\hertz} for every locality and ceilings configuration
        (\cref{sec:method}). The fixed clock makes the results reproducible, but it is not how the
        CPU normally runs. With turbo the compute-bound numbers scale up by roughly the clock ratio
        ($1.72\text{--}1.75\times$ in our turbo sweep) while DRAM bandwidth does not, so the
        DRAM-resident SAXPY speedups reported here ($1.88\text{--}2.02\times$) are larger than at turbo
        ($1.32\text{--}1.44\times$, \cref{sec:method}).
  \item \textbf{One machine, one compiler.} All results are from a single i9-10900F with GCC
        \sysGCC{}. Clang was used only to confirm the build compiles. No AVX-512 part was available,
        so the \qty{512}{bit} regime is out of scope.
  \item \textbf{THP always-on; no core isolation.} Transparent huge pages are \code{always}, and
        the pinned core is not isolated from the scheduler (no \code{isolcpus}). Neither is controlled
        for; the lock file only serializes cooperating jobs, so the shared-core checks of
        \cref{sec:method} (repetition spread $\le1.36\times$, apparent clock
        \qtyrange{2.78}{2.79}{\giga\hertz}) are what guard the reported data.
  \item \textbf{Discarded sessions.} Earlier sessions the same day were discarded: those before the
        GPE storm was masked (\cref{sec:method}), and a first clean session during which another
        agent's unlocked smoke tests briefly shared CPU~2 ($\approx2\times$ slowdowns in some
        repetitions).
  \item \textbf{Report build backend.} Without \code{biber}, \code{biblatex} uses the BibTeX
        backend, so the single reference is sorted by name/year to avoid a spurious rerun warning.
\end{itemize}
